\documentclass[book.tex]{subfiles}

\begin{document}

\chapter{Free Energy Perturbation}
\label{chap:fep}
\noindent\rule{11cm}{0.4pt} \\
\textbf{Prerequisites:} \autoref{chap:newtons_laws}, \autoref{chap:electrostatics}  \\
\textbf{Difficulty Level:} ** \\
\noindent\rule{11cm}{0.4pt}
\vspace{5mm}

Free energy perturbation techniques allow for estimates of the differences in free energy between different states of a physical system. Free energy perturbation systems can be critically important for computing quantities like the binding free energy of a molecule to a protein.

\section{Theoretical Foundations}

Suppose that we have two states of interest, state $0$ and state $1$. Let $A$ be the Helmholtz free energy. We can estimate the free energy state difference between state $0$ and $1$ by only drawing samples from state $0$ and then computing the potential energies for states $1$ and $0$, and then compute the estimated free energy difference using the update formula

\begin{align}
\Delta A &= A_1 - A_0 \\
&= -k_B T \ln \left \langle \exp \left (  -\frac{U_1 - U_0}{k_B T} \right ) \right \rangle
\end{align}
Here $U_1$ and $U_0$ are the potential energies of states $1$ and $0$ respectively. The challenge in using this equation though is that often state $1$ and state $0$ are far enough apart that computing the estimate directly from state $0$ can be prohibitively challenging. For this reason, the path from $0$ to $1$ is typically broken into a series of steps controlled by a parameter traditionally called $\lambda$.

\section{Targeted Approaches}

A new class of techniques seeks to make computations easier by learning a transformation that maps states $0$ and $1$ into overlapping regions. This learned transformation is typically a normalizing flow.
\end{document}

